\documentclass[conference]{IEEEtran}
\usepackage{amsmath,amssymb,graphicx,booktabs,siunitx,xcolor,multirow}
\usepackage[hidelinks]{hyperref}
\graphicspath{{../figures/out/}{fig/}}
% Figures are produced by run_all.py; a placeholder box is shown until they exist.
\newcommand{\fig}[2][\columnwidth]{\IfFileExists{../figures/out/#2.pdf}{\includegraphics[width=#1]{#2}}%
  {\IfFileExists{fig/#2.pdf}{\includegraphics[width=#1]{#2}}%
  {\fbox{\parbox[c][3cm][c]{\dimexpr#1-2\fboxsep-2\fboxrule}{\centering\small [\detokenize{#2}: produced by run\_all.py]}}}}}

\begin{document}
\title{Listen Before You Classify: A DSP Wake-Up Front End for Energy-Efficient and Reliable Sound Alerts on Wearables}
\author{\IEEEauthorblockN{Anonymous Authors}\IEEEauthorblockA{Anonymous Institution}}
\maketitle

\begin{abstract}
Wearable sound-awareness devices for Deaf and hard-of-hearing people run a neural audio classifier on
every half second of microphone input, which drains the battery, and pretrained classifiers raise
false alarms on safety classes. We place a classical signal-processing front end before the
classifier (YAMNet): a high-pass filter, short-time band energies and a per-band minimum-tracking
noise floor wake the classifier only when a new sound starts, and a Goertzel resonator bank with a
temporal-pattern matcher recognises standard alarm tones directly. The front end costs 0.5\,\% of one
classifier inference. On 8\,h of test streams built from real events and real backgrounds in six
environments never used for tuning, it lets the classifier sleep 86\,\% of the time while losing one of
288 detections, while fire-alarm recall on standard alarm tones rises from 29\,\% to 93\,\%. A diagnosis of the
classifier's false alarms (87\,\% are sirens labelled as fire alarms) leads to a veto that removes
83\,\% of them. On naturally recorded audio, the gate was awake for 97\,\% of 2\,765 labelled events
in 692 real home recordings and saves about 67\,\% of classifier compute over a real household's
day (99\,\% in an empty home, 23\,\% with the TV on), with no false alert in 5.5\,h. On held-out real
alarm recordings the DSP detector adds 5--13 points of smoke-alarm recall. We also show that
synthetic, standard-compliant alarms misjudge both detectors, and characterise real alarms to
explain why.
\end{abstract}

\begin{IEEEkeywords}
sound event detection, wake-up detector, Goertzel algorithm, noise-floor tracking, energy efficiency, assistive technology
\end{IEEEkeywords}

\section{Introduction}
Wearable sound-awareness devices turn important environmental sounds, such as smoke alarms, sirens and
door knocks, into vibration and visual alerts for Deaf and hard-of-hearing (DHH)
people~\cite{jain2020soundwatch}. They typically run a pretrained neural classifier such as
YAMNet~\cite{hershey2017cnn} on every half second of audio. This has two costs. \emph{Energy}: the
classifier runs continuously although nothing new happens most of the time, and battery drain was a key
constraint for on-watch sound recognition~\cite{jain2020soundwatch}. \emph{Reliability}: a general-purpose
classifier confuses related safety classes and can miss standardised signals such as the ISO~8201
temporal-three (T-3) smoke-alarm pattern~\cite{iso8201}.

Both problems have a classical signal-processing answer. Most of the day the acoustic scene is
stationary, which a cheap onset detector can tell without a neural network; and alarms are
\emph{designed} signals, with a known pitch and rhythm that a matched detector can recognise. We ask:
\emph{can an adaptive DSP front end that wakes the classifier only when a new sound starts, and that
detects alarm tones itself, cut classifier computation substantially without losing detections, in
acoustic environments it was never tuned on?}

Our contributions are:
\begin{enumerate}
\item A training-free front end (Butterworth high-pass, short-time band energies, minimum-tracking noise
floor, Goertzel resonator bank, T-3 matcher) costing 0.5\,\% of one classifier inference.
\item A controlled comparison against continuous inference, duty cycling and a fixed energy gate,
calibrated on 11 environments and tested on 6 unseen ones: 86\,\% less classifier compute with one of
288 detections lost (64\,\% with none lost), with a classifier duty that is far more stable across rooms than a fixed threshold.
\item A false-alarm diagnosis of the classifier and a one-line veto that removes 83\,\% of false alarms.
\item Validation on naturally recorded audio: 97\,\% of 2\,765 real labelled events wake the gate, the
compute saving over a real household's day is about 67\,\%, and held-out real alarm recordings show
where the DSP detector helps (in noise) and where it cannot (bells, horns, non-T-3 alarms).
\item Evidence that synthetic, standard-compliant alarms overstate a DSP detector (90\,\% vs.\ 25--30\,\%
on real recordings) and understate the classifier, with a characterisation of real alarms that
explains the gap. All data are public and the study is reproducible with one command.
\end{enumerate}

\section{Related Work}
\textbf{Sound awareness for DHH users.} SoundWatch~\cite{jain2020soundwatch} classified environmental
sounds on a smartwatch and compared watch-only, phone and cloud architectures, finding battery drain a key
concern; ProtoSound~\cite{jain2022protosound} personalised recognition from few examples. Neither reduces
\emph{when} the classifier runs.

\textbf{Wake-up detectors and cascades.} Always-on systems commonly place a low-power detector before a
costly stage~\cite{wakeupreview}; keyword spotting wakes large models with small
ones~\cite{gruenstein2017cascade}, and a recent cascade screens audio with a linear detector before a
trimmed YAMNet for violence detection~\cite{violencecascade2026}. We use interpretable, training-free DSP
gates whose threshold follows the room through noise-floor tracking, evaluate them on naturally recorded
home audio, and add a second, decision-making DSP path for alarm tones.

\textbf{Noise estimation and detection.} Our floor tracker is a simplified form of minimum
statistics~\cite{martin2001noise}; the ordered-statistic CFAR threshold we test as an ablation comes from
radar detection~\cite{rohling1983cfar}. Sigtia et al.\ compared environmental sound classifiers by accuracy
and computational cost, including smoke alarms recorded in homes~\cite{sigtia2016cost}.

\section{Method}
\subsection{Signals and classifier}
Audio is sampled at $f_s=16$\,kHz, the classifier's native rate. The Nyquist frequency of 8\,kHz covers the
piezo smoke-alarm band (about 2.8--3.6\,kHz) and the 520\,Hz low-frequency alarm~\cite{nfpa72} with its odd
harmonics up to the 13th. The classifier, YAMNet (MobileNetV1, 521 AudioSet classes), sees a 0.975\,s window
every 0.5\,s and costs $M_Y=\num{70744064}$ multiply-accumulates (MACs) per window, counted from its
convolution and dense layers. Its scores are mapped to three alert categories (fire alarm, siren, door) by
the maximum over category labels. An alert fires when the category score reaches 0.30 in $N$ classified
windows in a row ($N=2$ for fire alarm and siren, 1 for door; a run survives skipped windows if classified
windows are at most 2\,s apart), with a 20\,s cooldown and the most urgent category winning.

\subsection{Gate path}
\textbf{High-pass filter.} Butterworth, order 2, $f_c=100$\,Hz ($-3$\,dB), one biquad
\begin{equation}
H(z)=\frac{b_0+b_1z^{-1}+b_2z^{-2}}{1+a_1z^{-1}+a_2z^{-2}},
\end{equation}
with a double zero at $z=1$ and two poles inside the unit circle (Fig.~\ref{fig:filters}). It removes DC,
rumble and handling noise below the lowest target component, so they cannot wake the classifier.

\textbf{Band energies and noise floor.} Each 20\,ms frame ($N=320$, Hann) is transformed with a 512-point
FFT and its power summed in $B=16$ log-spaced bands (100\,Hz--7.6\,kHz), smoothed with a one-pole low-pass
($\tau=60$\,ms). With $E_b[t]$ the band level in dB, each band's floor follows the minimum with a bounded rise,
\begin{equation}
F_b[t]=\min\bigl(E_b[t],\,F_b[t-1]+\rho T\bigr),\qquad \rho=3\text{ dB/s},
\end{equation}
a simplified form of minimum statistics~\cite{martin2001noise}: it drops to a quieter background at once and
absorbs a steady sound within seconds.

\textbf{Wake-up.} The novelty $\nu[t]=\max_b(E_b[t]-F_b[t])$ is the largest rise of any band above its own floor,
so the threshold is relative to the room. Before each classifier window, if the maximum of $\nu$ over the
0.5\,s just received exceeds $\Delta$, the classifier runs, and stays awake for a hangover $H$. The window spans
the last 0.975\,s, so the onset that woke the gate lies inside the first classified window: waking adds no delay.

\subsection{Alarm path}
\textbf{Goertzel resonator bank.} The power of DFT bin $k$ in a frame is computed by the Goertzel
recursion~\cite{goertzel1958}
\begin{align}
s[n]&=x[n]+2\cos(\omega_k)\,s[n-1]-s[n-2],\\
|X_k|^2&=s_1^2+s_2^2-2\cos(\omega_k)\,s_1s_2,
\end{align}
a second-order IIR resonator with poles at $e^{\pm j\omega_k}$ on the unit circle, read once per frame. It
needs one multiply per sample per probe, against $N\log_2N$ for a full FFT, which pays off when few
frequencies matter. We use 17 probes from 2.8 to 3.6\,kHz (50\,Hz apart, one bin each) and 520, 1560 and 2600\,Hz.

\textbf{Tone track and pattern.} A frame is ``on'' when the strongest probe holds more than a share $r$ of the
frame energy (35\,\% for the three square-wave probes together) and is within 8\,dB of its 1.5\,s running
peak, which rejects reverberant tails. After closing 100\,ms dropouts, a fire alarm is declared on three
pulses of 0.3--0.8\,s separated by gaps of 0.3--0.8\,s, the T-3 timing with tolerance. The first design (v3)
used $r=25\,\%$ and also required the 1.5\,s pause after the third pulse (``strict''). The final design
(v4) was selected on real recordings (Sec.~\ref{sec:real}): $r=15\,\%$ without the pause check. Its alerts are
merged with the classifier's under the same cooldown.

\textbf{Siren veto.} A fire alarm from the classifier is accepted only when its score exceeds the siren score
in the same window (motivated in Sec.~\ref{sec:fa}).

\subsection{Front ends and cost}
F0 runs the classifier on every window; F1 on every $k$-th window; F2 wakes it when the high-passed level
exceeds a fixed $T$\,dBFS; F3 is the adaptive gate; F4, the proposed system, is F3 plus the alarm path and
the veto. Per 0.5\,s hop the gate costs \num{194535} MACs and the alarm path \num{177500}, counted per stage,
together 0.53\,\% of one classifier inference. Compute relative to F0 is $(d\,M_Y+M_\text{gate}+M_\text{alarm})/M_Y$,
with $d$ the classifier duty (fraction of windows classified).

\section{Experimental Setup}
\subsection{Test streams}
Following the event-mixing protocol of DCASE 2017 Task 2~\cite{mesaros2017dcase}, we build 5-minute streams
from DEMAND~\cite{thiemann2013demand} recordings (8 microphones per environment as separate streams, background
level $-50$ to $-35$\,dBFS). Each event stream holds 6 targets (2 per category) at an SNR from
$\{-5,0,5,10,20\}$\,dB over the event's active frames, and 6 distractors (0--20\,dB) that never overlap a
target; each has an \emph{idle} twin with the background only. Sirens, doorbells and knocks are real recordings
from FSD50K~\cite{fonseca2022fsd50k} and ESC-50~\cite{piczak2015esc} (official splits); distractors are other
FSD50K/ESC-50 sounds, half of them easily confused with targets (bells, phones, doors, impacts, beeps, speech,
music). Fire alarms are synthesised to ISO~8201 (piezo 2.9--3.5\,kHz or a 520\,Hz square wave, jittered timing
and pitch, a random room response), because no public set of real smoke-alarm recordings exists. Tuning uses
11 environments (7.3\,h of event and 7.3\,h of idle streams, 528 targets); all reported numbers use 6 others
(one per DEMAND category: living room, office, cafeteria, street, bus, park; 4.0\,h + 4.0\,h, 288 targets).

\subsection{Calibration}
For each gate we choose, on the tuning streams, the lowest-duty setting whose \emph{gate-caused} miss rate (the
share of F0's detections the gated system loses) stays within a budget of 1, 3 or 5\,\%. Four gate statistics
were compared on tuning data (Table~\ref{tab:ablation}); the selection rule kept the simplest.

\subsection{Real recordings}\label{sec:real}
\textbf{Alarms.} 161 Freesound~\cite{font2013freesound} recordings (CC0/CC-BY/CC-BY-NC): 38 smoke alarms, 30 fire
alarms (bells, horns, beeps) and 93 look-alikes that should not raise a fire alarm (microwave, oven and kitchen
timers, reversing beepers, alarm clocks, monitors, phones). They are split \emph{by uploader} into a design half and
a held-out half, so no recording device or room appears in both. The detector redesign searched probe band
(4 widths), tone share (3 values) and pattern rule (T-3, strict T-3, periodic beeps, combinations) on the design
half, with a rule fixed in advance: highest alarm recall with no fire alarm on design-half look-alikes. Each clip is
tested clean and mixed into a test-environment background at 10 and 0\,dB, after 5\,s of that background so the
noise floor has settled, as on an always-listening device.

\textbf{Real scenes.} DESED public evaluation set~\cite{turpault2019desed}: 692 real 10\,s home recordings with
time-stamped labels of 10 classes (2\,765 events). SINS~\cite{dekkers2017sins,dekkers2018dcase}: real recordings of one
home labelled with the household's activity; we score 1\,981 10\,s segments (5.5\,h, up to 400 per activity), each
after a 2\,s warm-up. Nothing is tuned on these data.

\subsection{Metrics}
A target is detected if an alert of its category fires between its onset and 1\,s after its end. We report recall
with 95\,\% bootstrap intervals that resample whole streams (events in one stream share a background), false alarms
per hour, median latency, classifier duty and relative compute; paired differences use the exact McNemar test.

\section{Results}
\subsection{Compute and recall}
Table~\ref{tab:main} gives the test results. Duty cycling (F1) cannot save anything: classifying every second window
already loses more than 5\,\% of detections on tuning data, so calibration returns $k=1$. Both gates keep F0's recall within one or two of 288 detections.
At the 5\,\% budget the adaptive gate F3 lets the classifier sleep 86\,\% of the time in background-only audio and
71\,\% in event streams, against 64\,\% and 52\,\% for the fixed gate F2, at the same recall; at the 3\,\% budget the
proposed system loses no detection at all with 64\,\% less compute. Fig.~\ref{fig:pareto} shows the full trade-off: the
adaptive gate dominates the fixed one at low duty, and duty cycling collapses.

\begin{table}[t]
\centering\caption{Test results (6 unseen environments, 288 targets). FA/h on event streams; duty in background-only
streams; compute relative to F0.}\label{tab:main}
\setlength{\tabcolsep}{3pt}
\begin{tabular}{lrrrrrrr}
\toprule
Front end & Recall & Fire & Siren & Door & FA/h & Duty & Comp. \\
 & (\%) & (\%) & (\%) & (\%) & & (\%) & \\
\midrule
F0 always on & 32.6 & 29.2 & 53.1 & 15.6 & 11.5 & 100 & 1.00 \\
F0 + veto & 32.6 & 29.2 & 53.1 & 15.6 & 2.0 & 100 & 1.00 \\
F0 + alarm DSP + veto & 53.8 & 92.7 & 53.1 & 15.6 & 2.0 & 100 & 1.00 \\
F1 every $k$-th & \multicolumn{7}{c}{no $k>1$ meets even the 5\,\% budget} \\
F2 fixed gate (5\,\%) & 31.9 & 28.1 & 53.1 & 14.6 & 10.5 & 36 & 0.36 \\
F3 adaptive gate (5\,\%) & 32.3 & 29.2 & 52.1 & 15.6 & 11.3 & 14 & 0.14 \\
F4 (1\,\% budget) & 53.8 & 92.7 & 53.1 & 15.6 & 2.0 & 66 & 0.67 \\
F4 (3\,\% budget) & 53.8 & 92.7 & 53.1 & 15.6 & 2.0 & 35 & 0.36 \\
\textbf{F4 proposed (5\,\%)} & \textbf{53.5} & \textbf{92.7} & 52.1 & 15.6 & \textbf{2.0} & \textbf{14} & \textbf{0.14} \\
F4 with DSP v3 (5\,\%) & 52.1 & 88.5 & 52.1 & 15.6 & 2.0 & 14 & 0.14 \\
\bottomrule
\end{tabular}
\end{table}

\begin{figure}[t]
\centering\fig{fig_pareto}
\caption{Recall against classifier duty on the test event streams for the duty-cycled (F1), fixed (F2) and adaptive
(F3) gates; markers: settings chosen on tuning data.}\label{fig:pareto}
\end{figure}

\subsection{Robustness across rooms}
A fixed level threshold depends on how loud a room is. In background-only audio its duty at the 5\,\% budget ranges
from 0.7\,\% (bus) to 69\,\% (cafeteria), while the adaptive gate stays between 1.4\,\% (street) and 24\,\% (bus) in all
six unseen rooms (Fig.~\ref{fig:idle}), because its threshold is relative to each band's own floor.

\begin{figure}[t]
\centering\fig{fig_idle_duty}
\caption{Classifier duty in background-only audio per unseen test environment.}\label{fig:idle}
\end{figure}

\textbf{Gate statistic ablation.} Longer smoothing with a two-band rule, and an ordered-statistic CFAR threshold
(median novelty of the last 10\,s)~\cite{rohling1983cfar}, did not improve on the simple statistic
(Table~\ref{tab:ablation}). CFAR was clearly worse: during a long sound such as a siren its reference median
rises, and the statistic sinks just when the classifier's scores peak.

\begin{table}[t]
\centering\caption{Classifier duty (\%) on the tuning environments within each gate-caused miss budget.}\label{tab:ablation}
\begin{tabular}{lrrr}
\toprule
Gate & 1\,\% & 3\,\% & 5\,\% \\
\midrule
F1 every $k$-th & 100 & 100 & 100 \\
F2 fixed level & 46.7 & 37.8 & 30.7 \\
F3 max band rise (selected) & 49.6 & 31.0 & 21.4 \\
\quad + 150\,ms smoothing, two bands & 50.6 & 28.6 & 23.1 \\
\quad + OS-CFAR & 85.2 & 47.1 & 36.1 \\
\quad + both & -- & 45.3 & 35.4 \\
\bottomrule
\end{tabular}
\end{table}

\subsection{False alarms}\label{sec:fa}
All of F0's false alarms occur in event streams (none in 4\,h of background): 40 of 46 are fire alarms raised during
sirens, because YAMNet also scores \emph{Alarm} on sirens and the fire alarm has priority. Tuning the classifier's
thresholds on the tuning streams (best 0.275 instead of 0.30) does not change recall or false alarms. The siren veto
cuts false alarms from 11.5 to 2.0 per hour (83\,\%) at unchanged recall and lowers the median alert latency from
1.85 to 1.55\,s. The DSP
alarm path adds no false alarm on the 1\,000+ everyday distractors.

\subsection{Alarm detection: synthetic versus real}
On the synthetic ISO~8201 alarms of the test streams, the DSP path raises fire-alarm recall from 29\,\% to 93\,\%
($p<10^{-15}$); the classifier never detects the 520\,Hz alarm, while the DSP path detects most of them (85\,\% in the first design). The gate does not
delay alerts (median 1.84\,s vs.\ 1.85\,s for F0); the alarm path needs about 2.4\,s to observe three pulses.

\begin{table}[t]
\centering\caption{Stress test (20 clips each, 5\,dB SNR): share raising a fire alarm.}\label{tab:stress}
\setlength{\tabcolsep}{3pt}
\begin{tabular}{lcrrr}
\toprule
Signal & Fire? & YAMNet & DSP strict & DSP v4 \\
\midrule
T-3, 3.1\,kHz piezo (in spec) & yes & 75 & 85 & 85 \\
T-3, 520\,Hz square (in spec) & yes & 0 & 90 & 90 \\
T-3, 2.6\,kHz / 4.0\,kHz & yes & 70 / 65 & 50 / 0 & 50 / 0 \\
T-3, timing $\times$0.7 / $\times$1.5 & yes & 90 / 70 & 35 / 35 & 35 / 30 \\
T-3, reverberation 1.5\,s & yes & 60 & 75 & 70 \\
T-4 carbon-monoxide pattern & no & 50 & 0 & 0 \\
Continuous 1\,Hz beeper, 3.1\,kHz & no & 80 & 0 & 85 \\
Single triple beep, 2.8\,kHz & no & 15 & 0 & 0 \\
\bottomrule
\end{tabular}
\end{table}

Real recordings tell a different story. Characterising the 32 design-half real alarms with DSP (FFT peak of tonal
frames, beep timing) shows a median pitch of 3.0\,kHz but only 54\,\% within the 2.8--3.6\,kHz probes (range
1.0--3.5\,kHz), and only 6 of 32 follow the T-3 pattern: 4 are continuous tones, 3 other rhythms, and 19 are not
narrowband tones at all (bells, horns, distorted recordings). A tone-and-pattern detector can therefore only reach a
minority of real alarms. The redesign raised design-half recall from 18\,\% to 25\,\% with no look-alike alarm;
wider bands or a periodic-beep rule reached 27--33\,\% but fired on 2--17\,\% of look-alikes and were rejected.

\begin{table}[t]
\centering\caption{Held-out real recordings (Freesound, unseen uploaders): share raising a fire alarm (\%), clean /
10\,dB / 0\,dB. F4 is the full gated system.}\label{tab:real}
\setlength{\tabcolsep}{3pt}
\begin{tabular}{lccc}
\toprule
Group ($n$) & YAMNet (F0) & DSP v4 alone & F4 proposed \\
\midrule
Smoke alarms (23) & 43 / 43 / 22 & 30 / 26 / 26 & \textbf{52 / 48 / 35} \\
Fire bells, horns (13) & 62 / 62 / 31 & 0 / 0 / 0 & 62 / 54 / 23 \\
Look-alikes (50) $\downarrow$ & 12 / 6 / 4 & 2 / 2 / 0 & 14 / 8 / 4 \\
\bottomrule
\end{tabular}
\end{table}

On held-out real recordings (Table~\ref{tab:real}) the classifier is much stronger than on synthetic alarms
(43--62\,\% clean vs.\ 29\,\%) and the DSP detector much weaker (26--30\,\% vs.\ 93\,\%). Combined, the proposed system
finds 5--13 points more smoke alarms than the always-on classifier, most at 0\,dB, where the classifier degrades and
the narrowband detector does not. The gate itself was awake during 98/91/78\,\% of the smoke alarms. It costs one of 13
fire bells at 10 and 0\,dB, where quiet bells do not always wake it. The DSP detector fired on 1 of 50 held-out
look-alikes, against 6 for the classifier.

\subsection{Naturally recorded audio}
On DESED's 692 real home recordings the gate was awake for 96.9\,\% of all labelled events (alarm/bell/ringing 99.0\,\%,
speech 98.9\,\%, dishes 95.1\,\%, vacuum cleaner 87.5\,\%, the lowest, as a steady noise is partly absorbed by the floor),
and the gated system kept 30 of the 32 alerts of the always-on classifier. Table~\ref{tab:sins} gives the classifier
duty over real home life. It follows the household: under 1\,\% in an empty home, 10\,\% while working, 77\,\% with the TV
on. Weighted by how SINS divides the day, the classifier runs 33\,\% of the time, a 67\,\% compute saving. Neither the
classifier (with the veto) nor the DSP path raised an alert in the 5.5\,h.

\begin{table}[t]
\centering\caption{Real home life (SINS, 5.5\,h): classifier duty of the proposed system per activity.}\label{tab:sins}
\setlength{\tabcolsep}{3pt}
\begin{tabular}{lrr|lrr}
\toprule
Activity & h & Duty & Activity & h & Duty \\
\midrule
Absence & 1.1 & 0.9\,\% & Eating & 0.3 & 43.6\,\% \\
Vacuuming & 0.1 & 7.4\,\% & Dishwashing & 0.2 & 50.6\,\% \\
Working & 1.1 & 10.2\,\% & Visitors, calls & 0.6 & 74.5\,\% \\
Cooking & 0.7 & 28.1\,\% & Watching TV & 1.1 & 76.7\,\% \\
\multicolumn{6}{c}{Weighted by the full SINS day: 33\,\% duty, 67\,\% saving} \\
\bottomrule
\end{tabular}
\end{table}

\begin{figure*}[t]
\centering\fig[0.95\textwidth]{fig_filters}
\caption{DSP design: high-pass filter magnitude response, pole-zero plot (HPF and one Goertzel resonator) and
Goertzel probe responses ($N=320$, Hann).}\label{fig:filters}
\end{figure*}

\section{Discussion}
\textbf{What the front end buys.} The main gain is compute with reliability: the same detections at a fraction of the
classifier's compute, fewer false alarms, and no added delay from waking. The saving depends on the scene and
should be reported per activity, not as one number: our mixed scenes suggest 86\,\%, a real household 67\,\%, an empty
home 99\,\%, a home with the TV on 23\,\%.

\textbf{Where DSP detection helps, and where it cannot.} A matched DSP detector is robust in noise and rarely fooled by
look-alikes, but it only reaches alarms that are the narrowband patterns it encodes. On real recordings that is a
minority, so it is a safety net beside the classifier, not a replacement. Synthetic, standard-compliant test alarms hide
this: they made the DSP detector look three times better and the classifier half as good as on real recordings.
Detector designs should be validated on real recordings, split by recording source.

\textbf{Limitations.} Real smoke-alarm recordings are scarce: 38 in total, so the held-out differences (2--3 clips) are
not statistically significant. SINS segments were scored independently after a short warm-up, not as continuous
episodes. We count multiply-accumulates, not battery energy; the microphone and the rest of the device keep their own
cost. The rule without the pause check (v4) fires on continuous 1\,Hz beepers (Table~\ref{tab:stress}); the strict rule
does not; the first, strict design reached 7 points less recall on the design-half real alarms, and the choice depends on the deployment. Fire bells and horns are
outside the detector's design. Results are for one classifier.

\section{Conclusion}
A classical front end costing 0.5\,\% of one neural inference lets a wearable's sound classifier sleep most of the time
without losing detections, adapts to unseen rooms where a fixed threshold does not, and, with a matched alarm detector and
a siren veto, raises alarm recall and removes most false alarms. Validation on naturally recorded audio confirmed the gate
and a 67\,\% compute saving over a real household's day, and showed that real alarms, unlike synthetic ones, are mostly
beyond a tone-and-pattern detector. Future work: continuous real-home recordings with labelled safety events, on-device
energy measurement, and detectors for the alarm types characterised here.

\bibliographystyle{IEEEtran}
\bibliography{refs}
\end{document}
